\documentclass[11pt,a4paper]{article}
\usepackage[margin=2.5cm]{geometry}
\usepackage{booktabs}
\usepackage{tabularx}
\usepackage{multirow}
\usepackage{graphicx}
\usepackage{xcolor}
\usepackage{hyperref}
\usepackage{enumitem}
\usepackage{amsmath}
\usepackage{float}

\definecolor{darkblue}{RGB}{0,51,102}
\hypersetup{colorlinks=true,linkcolor=darkblue,urlcolor=darkblue}

\title{\textbf{BGP-Sentry Dataset Properties\\[0.3em]\large Comprehensive Analysis and Realism Assessment}}
\author{BGP-Sentry-RS Evaluation Suite}
\date{April 2026}

\begin{document}
\maketitle
\tableofcontents
\newpage

% ============================================================================
\section{Dataset Overview}

BGP-Sentry uses four regional datasets extracted from real Internet topologies.
Each dataset represents a transit-layer backbone of one or more Regional Internet
Registry (RIR) service regions, preserving the natural provider--customer hierarchy,
peering density, and RPKI deployment rates observed in the wild.

\begin{table}[H]
\centering
\caption{Dataset Summary}
\label{tab:summary}
\begin{tabular}{lrrrrrr}
\toprule
\textbf{Dataset} & \textbf{ASes} & \textbf{RPKI} & \textbf{RPKI\%} & \textbf{Observers} & \textbf{Total Obs} & \textbf{Attacks} \\
\midrule
904\_afrinic\_transit\_mh   & 897   & 445   & 49.6\% & 897  & 222,908   & 20,373 \\
2030\_arin\_transit          & 2,020 & 1,371 & 67.9\% & 2,020 & 551,730  & 70,546 \\
3152\_lacnic\_afrinic\_transit & 3,104 & 1,930 & 62.2\% & 3,104 & 863,288 & 85,567 \\
5008\_lacnic\_5plus\_mh      & 4,997 & 2,671 & 53.5\% & 4,997 & 1,229,114 & 124,526 \\
\bottomrule
\end{tabular}
\end{table}

% ============================================================================
\section{Data Sources and Provenance}

\begin{table}[H]
\centering
\caption{Data Source Provenance}
\begin{tabular}{lll}
\toprule
\textbf{Component} & \textbf{Source} & \textbf{Type} \\
\midrule
AS identities            & CAIDA AS-relationships (2022-01-01) & Real \\
Customer--provider links & CAIDA (inferred from BGP feeds)      & Real \\
Peer links               & CAIDA (inferred from BGP feeds)      & Real \\
Regional AS allocation   & RIR delegation files                 & Real \\
RPKI signing status      & rpki-client VRP snapshot (2026-04-18)& Real \\
Per-AS prefixes          & pyasn (RouteViews RIB)               & Real \\
BGP observations         & BGPy Gao--Rexford simulation         & Simulated \\
Attack scenarios         & BGPy + \texttt{inject\_attacks.py}   & Synthetic \\
\bottomrule
\end{tabular}
\end{table}

The underlying CAIDA topology comprises 79,327 unique ASes, 163,285
customer--provider links, and 564,868 peer links. Regional datasets are
extracted by RIR allocation and transit hierarchy filtering.

% ============================================================================
\section{Topology Properties}

\subsection{Regional Extraction}

Each dataset is filtered by two criteria:
\begin{enumerate}
    \item \textbf{Region (RIR)} -- ASes allocated by a specific Regional Internet Registry.
    \item \textbf{Transit hierarchy} -- only ASes serving downstream customers are retained as validators.
\end{enumerate}

Stub ASes (leaf networks with no downstream customers) are excluded from the
validator set because they do not run RPKI validation independently, do not
participate in consensus, and BGP attacks propagate through the transit layer.

\subsection{RPKI Deployment Rates}

The RPKI percentages are \textbf{not artificially set}. They reflect real-world
signing status from the rpki-client VRP snapshot:

\begin{table}[H]
\centering
\caption{RPKI Deployment by Dataset}
\begin{tabular}{lrrrl}
\toprule
\textbf{Dataset} & \textbf{RPKI} & \textbf{Non-RPKI} & \textbf{RPKI\%} & \textbf{Notes} \\
\midrule
904 (AFRINIC)  & 445   & 459   & 49.6\% & Lower RPKI adoption in Africa \\
2030 (ARIN)    & 1,371 & 659   & 67.9\% & Higher North American adoption \\
3152 (LACNIC+AFRINIC) & 1,930 & 1,222 & 62.2\% & Mixed regional deployment \\
5008 (LACNIC 5+) & 2,671 & 2,337 & 53.5\% & Larger Latin American transit \\
\bottomrule
\end{tabular}
\end{table}

\subsection{Connectivity}

All datasets have 98\%+ of nodes in a single connected component. Isolated nodes
(degree 0 after extraction) are dropped to ensure every participant can reach
every other participant through the topology.

% ============================================================================
\section{Observation Properties}

\subsection{Per-Node Distribution}

\begin{table}[H]
\centering
\caption{Observations per Node}
\begin{tabular}{lrrrr}
\toprule
\textbf{Dataset} & \textbf{Min} & \textbf{Median} & \textbf{Max} & \textbf{Mean} \\
\midrule
904   & 2   & 198 & 2,233 & 248.5 \\
2030  & 102 & 222 & 1,538 & 273.1 \\
3152  & 1   & 243 & 8,965 & 278.1 \\
5008  & 1   & 207 & 5,864 & 246.0 \\
\bottomrule
\end{tabular}
\end{table}

The heavy-tailed distribution (max $\gg$ mean) reflects the real Internet topology
where highly-connected transit providers see many more BGP announcements than
edge ASes. For example, in the 3152 dataset, the busiest transit AS observes
8,965 announcements while the median AS sees only 243---a 37$\times$ ratio that
mirrors the real-world power-law distribution of BGP table sizes.

\subsection{Temporal Properties and Real-World Calibration}

\begin{table}[H]
\centering
\caption{Temporal Span and Event Rates}
\begin{tabular}{lrrll}
\toprule
\textbf{Dataset} & \textbf{BGP Span} & \textbf{Rate/Validator} & \textbf{RIPE RIS Target} & \textbf{Ratio} \\
\midrule
904   & 10h 54m & 0.0073 ev/s & 0.005 ev/s & 1.46$\times$ \\
2030  & 11h  9m & 0.0073 ev/s & 0.005 ev/s & 1.46$\times$ \\
3152  & 11h 50m & 0.0068 ev/s & 0.005 ev/s & 1.36$\times$ \\
5008  & 10h 38m & 0.0065 ev/s & 0.005 ev/s & 1.30$\times$ \\
\bottomrule
\end{tabular}
\end{table}

The per-validator event rate is calibrated against RIPE RIS 2022 steady-state
measurements ($\sim$0.005 events/sec/validator). The observed overshoot of
1.30--1.46$\times$ arises because BGPy's Gao--Rexford propagation generates
marginally more observation events per node than real BGP feeds. This is
considered acceptable for evaluation---within the same order of magnitude.

\subsection{Simulation Timing and Speed Scaling}

The simulation replays BGP timestamps in accelerated wall-clock time:

\begin{table}[H]
\centering
\caption{Simulation Timing at 20$\times$ Speed}
\begin{tabular}{lrrl}
\toprule
\textbf{Parameter} & \textbf{Base Value} & \textbf{Scaled (20$\times$)} & \textbf{Purpose} \\
\midrule
\texttt{RPKI\_DEDUP\_WINDOW}       & 15s   & 1s   & Skip repeated (prefix, origin) \\
\texttt{VOTING\_OBSERVATION\_WINDOW} & 300s & 15s  & KB entry validity for voting \\
\texttt{FLAP\_WINDOW\_SECONDS}     & 60s   & 3s   & Flap counting window \\
\texttt{FLAP\_DEDUP\_SECONDS}      & 30s   & 2s   & Flap event deduplication \\
\texttt{P2P\_REGULAR\_TIMEOUT}     & 30s   & 30s  & Consensus timeout (wall-clock) \\
\texttt{P2P\_ATTACK\_TIMEOUT}      & 60s   & 60s  & Attack TX timeout (wall-clock) \\
\bottomrule
\end{tabular}
\end{table}

\textbf{Key design decision}: BGP-time-sensitive parameters (dedup, KB window,
flap window) are divided by the speed multiplier so that their effective
BGP-time coverage remains unchanged. P2P consensus timeouts are \emph{not}
scaled because they govern wall-clock waiting for vote responses---they need
real seconds regardless of simulation speed.

\subsection{Relationship Distribution}

\begin{table}[H]
\centering
\caption{Receiver Relationship Distribution (\%)}
\begin{tabular}{lrrrr}
\toprule
\textbf{Dataset} & \textbf{CUSTOMERS} & \textbf{PROVIDERS} & \textbf{PEERS} & \textbf{ORIGIN} \\
\midrule
904   & 44.0\% & 24.2\% & 31.1\% & 0.7\% \\
2030  & 43.4\% & 25.1\% & 31.0\% & 0.4\% \\
3152  & 35.0\% & 25.9\% & 38.9\% & 0.2\% \\
5008  & 28.3\% & 25.6\% & 46.0\% & 0.1\% \\
\bottomrule
\end{tabular}
\end{table}

\textbf{Trend}: As topology size increases, the peering ratio grows (31\% $\to$
46\%) while customer-learned routes decrease (44\% $\to$ 28\%). This is
consistent with larger networks having denser IXP participation. The provider
ratio remains stable at $\sim$25\% across all sizes.

\subsection{Hop Distance Distribution}

\begin{table}[H]
\centering
\caption{AS-Path Hop Distance Distribution (\%)}
\begin{tabular}{lrrr}
\toprule
\textbf{Dataset} & \textbf{Hop 0 (self)} & \textbf{Hop 1 (neighbor)} & \textbf{Hop 2} \\
\midrule
904   & 95.9\% & 1.8\% & 2.4\% \\
2030  & 99.4\% & 0.6\% & 0.04\% \\
3152  & 97.7\% & 1.5\% & 0.8\% \\
5008  & 98.2\% & 0.6\% & 1.1\% \\
\bottomrule
\end{tabular}
\end{table}

The vast majority of observations are at hop distance 0 (self-origin), which are
filtered by the trusted-path filter. This means the hop limit setting dramatically
affects how many observations reach the detection pipeline---changing from
hop=1 to hop=2 can increase processed observations by 5--50$\times$.

% ============================================================================
\section{Attack Properties}

\subsection{Attack Type Distribution}

\begin{table}[H]
\centering
\caption{Attack Observations by Type and Category}
\begin{tabular}{llrrrr}
\toprule
\textbf{Category} & \textbf{Attack Type} & \textbf{904} & \textbf{2030} & \textbf{3152} & \textbf{5008} \\
\midrule
Instability     & ROUTE\_FLAPPING      & 10,629 & 47,406 & 48,929 & 69,633 \\
Invalid Inject. & BOGON\_INJECTION     & 4,728  & 9,996  & 17,298 & 24,473 \\
Prefix Hijack   & SUBPREFIX\_HIJACK    & 4,077  & 9,886  & 19,070 & 25,005 \\
Path Manip.     & PATH\_POISONING      & 926    & 3,219  & 85     & 5,414 \\
Policy Viol.    & ROUTE\_LEAK          & 13     & 39     & 185    & 1 \\
\midrule
                & \textbf{Total}       & 20,373 & 70,546 & 85,567 & 124,526 \\
\bottomrule
\end{tabular}
\end{table}

\subsection{Attack Rate and Realistic Context}

The overall attack observation ratio ranges from 9.1\% to 12.8\%. This is
deliberately higher than real-world rates ($<$0.1\% of BGP announcements):

\begin{itemize}
    \item Each unique attack event ($\sim$5 per type from BGPy) is observed by many
          validators through BGP propagation, amplifying observation counts.
    \item The \textbf{event-level} attack density is $\sim$1--2\%, much closer to
          realistic rates.
    \item In real BGP, a single prefix hijack event generates thousands of
          UPDATE messages across the DFZ---the same amplification seen here.
\end{itemize}

\subsection{Example: How an Attack Propagates and Gets Detected}

Consider a BOGON\_INJECTION attack where AS\,666 announces the reserved prefix
\texttt{10.0.0.0/8}:

\begin{enumerate}
    \item AS\,666 injects the bogon announcement into the BGP topology.
    \item The announcement propagates through Gao--Rexford valley-free routing
          to AS\,666's providers, who propagate it to their peers and customers.
    \item Each RPKI validator that receives the announcement (within its hop
          radius) checks it against the 13 reserved/private IP ranges.
    \item \texttt{10.0.0.0/8} matches the RFC\,1918 bogon range---the detector
          flags it as BOGON\_INJECTION with CRITICAL severity.
    \item The validator creates a signed transaction recording the detection
          and broadcasts it to peers for consensus voting.
    \item Peers who also observed the same (prefix, origin) vote APPROVE from
          their knowledge bases.
    \item With $\geq$3 APPROVE votes, the transaction is committed as CONFIRMED
          to each validator's independent blockchain.
\end{enumerate}

\subsection{Attack Generation Methods}

\begin{itemize}
    \item \textbf{SUBPREFIX\_HIJACK}, \textbf{BOGON\_INJECTION}, \textbf{ROUTE\_FLAPPING}:
          Generated by BGPy's built-in attack scenarios over the Gao--Rexford model.
    \item \textbf{ROUTE\_LEAK}, \textbf{PATH\_POISONING}: Injected by
          \texttt{scripts/inject\_attacks.py} because BGPy's default scenarios do not
          exercise these two vectors. Both achieve 100\% recall on injected events and
          0\% false positives on 300 legitimate samples per dataset.
    \item \textbf{ACCIDENTAL\_ROUTE\_LEAK}: Filtered out from datasets (its
          forged-origin pattern does not match the valley-free ROUTE\_LEAK detector).
\end{itemize}

% ============================================================================
\section{Realism Assessment}

\subsection{What is Realistic}

\begin{enumerate}
    \item \textbf{Topology}: Extracted from real CAIDA AS-relationship data, preserving
          natural provider--customer hierarchy and peering density.
    \item \textbf{RPKI deployment}: Real-world signing status from rpki-client VRP
          snapshots, not artificial percentages.
    \item \textbf{Event rate}: Per-validator observation rate (0.005--0.007 events/sec)
          closely matches RIPE RIS 2022 measurements.
    \item \textbf{Routing model}: BGPy implements Gao--Rexford valley-free routing,
          the standard model for BGP simulation in academic literature.
    \item \textbf{Propagation}: Announcements propagate through the topology naturally,
          producing multi-observer views of the same event.
    \item \textbf{Regional diversity}: Four datasets spanning AFRINIC, ARIN, and LACNIC
          regions with different RPKI adoption rates and peering densities.
\end{enumerate}

\subsection{Known Limitations}

\begin{enumerate}
    \item \textbf{Simulated BGP, not captured}: Observations come from BGPy simulation,
          not live BGP feeds (RIPE RIS, RouteViews). This means no BGP convergence
          delays, no MRAI timers, no route dampening, and no traffic engineering
          artifacts.
    \item \textbf{Static topology}: The CAIDA snapshot is from January 2022.
    \item \textbf{Small unique event count}: $\sim$5 unique attack events per type.
    \item \textbf{Uniform attack timing}: Attacks are injected uniformly; real attacks
          cluster temporally.
    \item \textbf{Limited prefix diversity}: Only 90--96 unique prefixes vs.\ 900,000+
          in real BGP tables.
    \item \textbf{Partial ROA coverage}: The 730,808-entry ROA database covers
          $\sim$46\% of simulated sub-prefix hijack attacks.
\end{enumerate}

\subsection{Comparison with Related Work}

\begin{table}[H]
\centering
\caption{Comparison with BGP Security Evaluation Datasets}
\begin{tabularx}{\textwidth}{lXXX}
\toprule
\textbf{Property} & \textbf{BGP-Sentry} & \textbf{BGPStream/RIPE RIS} & \textbf{ARTEMIS} \\
\midrule
Topology source   & CAIDA real         & Live Internet    & Live Internet \\
BGP model         & BGPy Gao--Rexford  & Real BGP feeds   & Real BGP feeds \\
Attack injection  & Synthetic          & Natural + staged & Natural \\
Scale             & 897--4,997 ASes    & Full Internet    & Per-prefix \\
Reproducibility   & Full               & Partial          & Low \\
Ground truth      & Complete labels    & Partial (manual) & Per-prefix \\
\bottomrule
\end{tabularx}
\end{table}

% ============================================================================
\section{Discussion Summary}

Our evaluation uses four regional datasets extracted from real CAIDA AS-relationship data (January 2022), ranging from 897 to 4,997 ASes across AFRINIC, ARIN, and LACNIC regions. The RPKI deployment rates (49.6\%--67.9\%) are not artificially set---they reflect actual rpki-client VRP snapshots. The per-validator event rate of 0.0065--0.0073 events/sec closely matches the real-world steady-state of $\sim$0.005 events/sec measured by RIPE RIS in 2022, only about 1.3--1.5$\times$ higher. BGP observations are generated by BGPy's Gao--Rexford model, which is the standard simulation model in BGP security literature. The main limitations are: observations are simulated rather than captured from live BGP feeds (so we miss convergence delays, MRAI timers, and route dampening), the topology is a static snapshot, and we only have $\sim$90 unique prefixes compared to 900,000+ in real BGP tables. Despite these limitations, the datasets provide reproducible, fully-labeled ground truth---something impossible with live BGP data---built on real topology and real RPKI deployment patterns.

\end{document}
